\begin{summarybox}
	\textbf{\textcolor{CSummary}{一句话核心}}
	本文采用“试验次数型几何分布”：若 $\xi$ 表示首次成功所需的总试验次数，则
	\[
		E\xi=\frac{1}{p}.
	\]

	\medskip

	\textbf{\textcolor{CSummary}{使用条件}}

	\begin{itemize}[leftmargin=*, itemsep=0.5em, topsep=0.4em]
		\item \textbf{条件 1（独立重复）：} 各次试验相互独立。
		\item \textbf{条件 2（概率恒定）：} 每次试验成功概率均为 $p$，失败概率为 $q=1-p$。
		\item \textbf{条件 3（定义匹配）：} 随机变量 $\xi$ 表示首次成功所需的总试验次数，因此 $\xi=1,2,\dots$。
	\end{itemize}

	\medskip

	\textbf{\textcolor{CSummary}{关键公式}}

	\[
		P(\xi=k)=q^{k-1}p,\quad k=1,2,\dots,\qquad E\xi=\frac{1}{p}.
	\]

	\medskip

	\textbf{\textcolor{CSummary}{定义辨析}}

	\begin{itemize}[leftmargin=*, itemsep=0.5em, topsep=0.4em]
		\item \textbf{总试验次数型：} 若 $\xi$ 表示首次成功所需总试验次数，则 $E\xi=\dfrac{1}{p}$。
		\item \textbf{失败次数型：} 若 $X$ 表示首次成功前的失败次数，则 $X=\xi-1$，所以 $EX=\dfrac{q}{p}$。
	\end{itemize}

	\medskip

	\textbf{\textcolor{CSummary}{关键提醒}}

	\begin{itemize}[leftmargin=*, itemsep=0.5em, topsep=0.4em]
		\item \textbf{易错点（定义混淆）：} 先看随机变量表示“总试验次数”还是“失败次数”，再决定使用 $\dfrac{1}{p}$ 还是 $\dfrac{q}{p}$。
		\item \textbf{易错点（pq混淆）：} 本文公式分母是成功概率 $p$，不是失败概率 $q$。
		\item \textbf{检查项（参数范围）：} 使用前确认 $0<p\le1$；反解参数时还要检查 $p$ 是否落在该范围内。
		\item \textbf{检查项（模型前提）：} 若各次试验不独立，或成功概率会变化，则不能直接套用普通几何分布公式。
	\end{itemize}
\end{summarybox}